\section{Conclusion}
\label{sec:conclusion}

Courts that compel governments to buy medications save individual lives. They also cost the public budget \BPnegPctHeadline{} more per purchase across the full sample, attract \BPfirmsPctHeadline{} fewer bidders, and---when the procurement officer is exposed to personal sanctions---generate a \BPutgPctRange{} cost margin between litigated and administrative urgent purchases of the same item.

The decomposition we develop traces this margin to three distinct channels operating simultaneously. \emph{Demand fragmentation} (C1) does most of the work: sanctions force smaller orders, and the bulk-discount channel mechanically explains essentially all of the under-the-gun gap and roughly half of the urgency premium across the full sample. \emph{Reduced participation} (C2) accounts for most of the across-sample residual: at identical order size, urgent tenders attract \BPfirmsPanelBpct{} fewer bidders, leaving thinner markets and worse negotiated prices. \emph{Supplier composition shift} (C3) accounts for the rest of the across-sample premium: under urgency, the supplier mix shifts toward firms that specialize in compressed-timeline delivery, yielding worse matches on average without any within-firm markup. The within firm-buyer-item triple test directly rules out the within-firm-markup interpretation that supplier fixed effects might suggest in isolation: the same firm sells the same item to the same buyer at a slightly \emph{lower} per-unit price under urgency, not a higher one.

The decomposition has direct policy content. Two complementary reforms address two distinct channels.

On the demand side, expanding S\~ao Paulo's administrative-request mechanism---which since \BPsampleStartYear{} has allowed urgent procurement without personal-sanction exposure---would eliminate the channel that operates inside UTG. Applied to a defensible target of \BPpolicyAdminEligibleShare{} of court-mandated spending, this would recover on the order of \BPpolicyAdminRecoveryHigh{} per year on the \BPlitSpendingUSDmm~M of S\~ao Paulo's annual litigated spending. The instrument is administrative capacity (the scientific-committee infrastructure), not legal reform.

On the supply side, framework agreements for items where fragmentation is the binding cost---repeat-purchase commodity-like medications---directly address C1 by allowing aggregation across urgent and ordinary requests. Calibrated to a feasible \BPpolicyFrameworkFeasibleShare{} of the annual spending, the C1 channel mapped onto this scope would recover on the order of \BPpolicyFrameworkRecovery{} per year. Limiting the prominence of court-order references in tender notices (which currently serves as a public signal of the government's reduced outside option) would address C3 without constraining the underlying litigation; we view this as second-order relative to the C1 reform.

The fiscal scale of the cost is substantial but bounded. Applying the urgent-vs-ordinary premium to the \$\BPlitSpendingUSDmm~M of annual litigated spending in S\~ao Paulo gives a per-unit-price-margin lower bound of \BPwelfareBoundLow{} per year. Applying the under-the-gun premium to the \$\BPsanctionExposedUSDmmLow--\BPsanctionExposedUSDmmHigh~M of sanction-exposed spending gives an upper bound of \BPwelfareBoundHigh{} per year. The publishable range is \BPwelfareBoundRange{} per year on the per-unit-price margin alone; adding the welfare cost of smaller order sizes, reduced competition, and officials' search time diverted from ordinary procurement would raise this further.\footnote{The lower bound applies the headline urgent-vs-ordinary premium (\BPnegPctHeadline) to the entire \$\BPlitSpendingUSDmm~M base. The upper bound applies the under-the-gun premium (\BPutgPctRange) to the sanction-exposed portion (\$\BPsanctionExposedUSDmmLow--\BPsanctionExposedUSDmmHigh~M). Both are per-unit-price-margin estimates and exclude welfare losses from reduced competition and constrained order sizes.}

These findings highlight a tension at the heart of right-to-health enforcement: judicial mandates intended to guarantee individual access impose substantial costs on the public budget, potentially reducing resources available for the broader health system. Institutional designs that balance individual rights with procurement efficiency---rather than treating them as competing objectives---could improve welfare for both litigants and the broader population. Brazil's recent procurement reform (Lei~14.133/2021), which expands electronic auctions and introduces framework agreements, could mitigate the quantity-fragmentation channel we document---a prediction future work can test as the reform takes effect.
